\documentclass[../../../include/open-logic-section]{subfiles}

\begin{document}

\olfileid{sth}{replacement}{intro}
\olsection{مقدمة}

الاستبدال هو مخطط المسلّمات الذي يصنع الفارق بين~$\ZF$ و$\Z$. وقد
استعملناه من دون إثبات طوال
\crefrange{sth:ordinals::chap}{sth:spine::chap}. وسنبحث أخيرًا في هذا
الفصل السؤال: هل الاستبدال مسوَّغ؟

ولجعل السؤال محددًا، يجدر أن نلاحظ أن الاستبدال \emph{قوي} حقًا.
وسندرك في هذا الفصل مقدار قوته بالضبط (ثم ندركه مرة أخرى في
\olref[sth][card-arithmetic][fix]{sec}). لكن ذلك سيوحي بأن التسويغ
مطلوب حقًا.

سنناقش نوعين من التسويغ. وعلى وجه التقريب: يحاول التسويغ
\emph{الخارجي} تسويغ مسلّمة بثمارها؛ ويحاول التسويغ \emph{الداخلي}
تسويغها باقتراح أن المفاهيم الرياضية المعنية تصدقها. وسيتضح معنى هذا
أكثر في أثناء الفصل، لكنه ليس إلا رأس جبل جليدي. وللمزيد، انظر خصوصًا
\citeauthor{Maddy1988a} (\citeyear{Maddy1988a} و
\citeyear{Maddy1988b}).

\end{document}
